\section{Terminologie et concepts de base}
\addcontentsline{toc}{section}{Terminologie et concepts clés}

Afin de faciliter la lecture de ce document et la compréhension des enjeux techniques liés à la transmission 
vidéo basse latence, les concepts et acronymes suivants seront couramment utilisés :

\begin{description}
    \item[Wi-Fi HaLow (IEEE 802.11ah) :] Amendement du standard Wi-Fi conçu pour opérer dans les bandes de 
    fréquences inférieures à 1 GHz (Sub-1 GHz). Contrairement au Wi-Fi classique (2.4 GHz ou 5.8 GHz), il 
    privilégie la portée (jusqu'à plusieurs kilomètres) et la capacité à traverser les obstacles physiques, 
    au détriment de la bande passante maximale.
    
    \item[FPV (\textit{First Person View}) :] "Vol en immersion". Technique de pilotage où la caméra embarquée 
    sur le drone transmet un flux vidéo en temps réel vers les lunettes ou l'écran du pilote, exigeant une 
    latence extrêmement faible et constante.
    
    \item[NLOS (\textit{Non-Line-Of-Sight}) :] Situation de communication radio où l'émetteur et le récepteur 
    ne sont pas en ligne de vue directe (présence de bâtiments, de végétation ou de relief). C'est dans ces 
    conditions que les hautes fréquences perdent leur efficacité.
    
    \item[MCS (\textit{Modulation and Coding Scheme}) :] Indice (généralement de 0 à 10) définissant le débit 
    physique de la transmission. Il combine un type de modulation radio (ex: BPSK, QAM) et un taux de codage. 
    Un MCS bas (ex: 1 ou 2) offre un débit réduit, mais garantit une robustesse maximale face au bruit 
    et aux interférences.
    
    \item[LDPC (\textit{Low-Density Parity-Check}) :] Algorithme avancé de correction d'erreurs (FEC). Contrairement 
    au codage convolutif standard (BCC), le LDPC utilise des matrices de parité complexes permettant au récepteur de 
    reconstruire des paquets de données corrompus par de mauvaises conditions radio, évitant ainsi de devoir redemander 
    l'envoi du paquet (ce qui détruirait la latence vidéo).
    
    \item[Couches PHY et MAC :] Niveaux fondamentaux du modèle réseau. La \textbf{couche PHY} (Physique) gère la transformation 
    des bits en ondes radio. La \textbf{couche MAC} (\textit{Media Access Control}) gère les règles d'accès à l'antenne et 
    les retransmissions. Ce projet modifie directement ces couches pour les adapter au flux vidéo.
    
    \item[RTOS (\textit{Real-Time Operating System}) :] Système d'exploitation temps réel (comme FreeRTOS, utilisé sur 
    le microcontrôleur d'émission STM32). Contrairement à un OS classique (comme Linux), un RTOS garantit l'exécution des tâches dans 
    des délais stricts et déterministes, minimisant ainsi la gigue logicielle lors du routage des données.

    \item[Latence "Glass-to-Glass" :] Délai total mesuré entre l'instant où un événement physique se produit devant l'objectif 
    de la caméra (le premier "verre") et le moment où il est affiché sur l'écran du pilote (le second "verre"). C'est la métrique 
    de performance absolue de ce projet, englobant tous les temps de traitement matériels et logiciels.

    \item[UDP Broadcast \& No-ACK :] Stratégie de transmission réseau consistant à diffuser des paquets de données à l'aveugle à tous 
    les équipements du sous-réseau, sans demander d'accusé de réception (\textit{No-ACK}). Cette méthode sacrifie la garantie absolue de 
    livraison au profit d'un flux continu à latence minimale.

    \item[Gigue (\textit{Jitter}) :] Variation temporelle de la latence. Dans un flux vidéo en temps réel, une forte gigue (des paquets arrivant 
    en rafales irrégulières) provoque des saccades visuelles souvent plus perturbantes pour le pilotage qu'une latence fixe et stable.
\end{description}

\vspace{1cm}